\textbf{Phần 1. Trắc nghiệm lựa chọn}
\setcounterpageref{ex}{0}

\begin{ex} %[3D2V29]
Tính giá trị của biểu thức: $M = 84 : 4 \times 3 - 25$.
\choice
{\True $38$}
{$42$}
{$35$}
{$40$}
\loigiai{
$84 : 4 = 21$; $21 \times 3 = 63$; $63 - 25 = 38$.
\textbf{Đáp án của câu hỏi là A}
}
\end{ex}

\begin{ex} %[3D2V29]
Tìm $x$, biết: $x \times 5 + 16 = 86$.
\choice
{$12$}
{$16$}
{\True $14$}
{$15$}
\loigiai{
$x \times 5 = 86 - 16 = 70 \Rightarrow x = 70 : 5 = 14$.
\textbf{Đáp án của câu hỏi là C}
}
\end{ex}

\begin{ex} %[3D2V29]
Một cửa hàng có 96 chiếc áo sơ mi. Ngày thứ nhất bán được $\dfrac{1}{4}$ số áo đó. Ngày thứ hai bán được $\dfrac{1}{3}$ số áo còn lại. Hỏi sau hai ngày cửa hàng còn lại bao nhiêu chiếc áo sơ mi?
\choice
{$54$ chiếc áo}
{\True $48$ chiếc áo}
{$40$ chiếc áo}
{$36$ chiếc áo}
\loigiai{
Ngày thứ nhất bán: $96 : 4 = 24$ chiếc áo. Còn lại sau ngày 1: $96 - 24 = 72$ chiếc áo.\
Ngày thứ hai bán: $72 : 3 = 24$ chiếc áo. Sau 2 ngày còn lại: $72 - 24 = 48$ chiếc áo.
\textbf{Đáp án của câu hỏi là B}
}
\end{ex}

\begin{ex} %[3D2V29]
Trong một phép chia có số chia là 6, thương là 15 và số dư là số dư lớn nhất có thể. Số bị chia là
\choice
{$90$}
{$94$}
{$96$}
{\True $95$}
\loigiai{
Số dư lớn nhất là 5. Số bị chia là: $6 \times 15 + 5 = 90 + 5 = 95$.
\textbf{Đáp án của câu hỏi là D}
}
\end{ex}

\begin{ex} %[3D2V29]
Một khu vườn hình chữ nhật có chu vi $72\text{ m}$, chiều rộng là $14\text{ m}$. Chiều dài của khu vườn đó gấp mấy lần chiều rộng?
\choice
{\True Không chia hết, chiều dài là $22\text{ m}$}
{Gấp 2 lần}
{Gấp 3 lần}
{Gấp 4 lần}
\loigiai{
Nửa chu vi: $72 : 2 = 36\text{ m}$. Chiều dài: $36 - 14 = 22\text{ m}$. Chiều dài không chia hết cho chiều rộng ($22 : 14$).
\textbf{Đáp án của câu hỏi là A}
}
\end{ex}

\begin{ex} %[3D2V29]
Cho dãy số: $3, 6, 12, 24, \dots$ Quy luật của dãy số là số đứng sau gấp 2 lần số đứng liền trước. Số hạng tiếp theo của dãy số là
\choice
{$36$}
{$40$}
{\True $48$}
{$50$}
\loigiai{
Số hạng tiếp theo là: $24 \times 2 = 48$.
\textbf{Đáp án của câu hỏi là C}
}
\end{ex}

\begin{ex} %[3D2V29]
Có 4 hộp bút, mỗi hộp có 12 chiếc bút. Người ta lấy ra 18 chiếc bút để chia cho các bạn học sinh giỏi. Số bút còn lại được chia đều vào 6 túi. Mỗi túi có bao nhiêu chiếc bút?
\choice
{$6$ chiếc bút}
{\True $5$ chiếc bút}
{$4$ chiếc bút}
{$7$ chiếc bút}
\loigiai{
Tổng số bút: $12 \times 4 = 48$ chiếc. Số bút còn lại: $48 - 18 = 30$ chiếc. Mỗi túi có: $30 : 6 = 5$ chiếc bút.
\textbf{Đáp án của câu hỏi là B}
}
\end{ex}

\begin{ex} %[3D2V29]
Nếu lấy một số nhân với 4 rồi trừ đi 15 thì được 65. Số đó là
\choice
{$18$}
{$22$}
{$25$}
{\True $20$}
\loigiai{
Số đó nhân 4 là: $65 + 15 = 80$. Số cần tìm là: $80 : 4 = 20$.
\textbf{Đáp án của câu hỏi là D}
}
\end{ex}

\begin{ex} %[3D2V29]
Bạn Cúc gấp được 18 ngôi sao, bạn Lan gấp được số ngôi sao gấp đôi bạn Cúc. Hỏi cả hai bạn gấp được tất cả bao nhiêu ngôi sao?
\choice
{\True $54$ ngôi sao}
{$36$ ngôi sao}
{$48$ ngôi sao}
{$60$ ngôi sao}
\loigiai{
Lan gấp được: $18 \times 2 = 36$ ngôi sao. Cả hai bạn gấp được: $18 + 36 = 54$ ngôi sao.
\textbf{Đáp án của câu hỏi là A}
}
\end{ex}

\begin{ex} %[3D2V29]
Thực hiện phép chia: $89 : 7$. Kết quả của phép chia là
\choice
{Thương $12$, dư $3$}
{Thương $11$, dư $12$}
{\True Thương $12$, dư $5$}
{Thương $13$, dư $2$}
\loigiai{
Ta có: $89 = 7 \times 12 + 5$ với số dư $5 < 7$.
\textbf{Đáp án của câu hỏi là C}
}
\end{ex}

\begin{ex} %[3D2V29]
Một đoàn khách có 43 người cần thuê xe ô tô 7 chỗ (không kể tài xế). Cần thuê ít nhất bao nhiêu xe để chở hết đoàn khách đó?
\choice
{$6$ xe}
{\True $7$ xe}
{$8$ xe}
{$5$ xe}
\loigiai{
Ta có: $43 : 7 = 6$ (dư 1 người). Cần ít nhất $6 + 1 = 7$ xe ô tô.
\textbf{Đáp án của câu hỏi là B}
}
\end{ex}

\begin{ex} %[3D2V29]
Tổng của hai số là 64. Biết số lớn gấp 3 lần số bé. Số bé là
\choice
{$12$}
{$14$}
{$18$}
{\True $16$}
\loigiai{
Tổng số phần bằng nhau: $3 + 1 = 4$ phần. Số bé là: $64 : 4 = 16$.
\textbf{Đáp án của câu hỏi là D}
}
\end{ex}

\textbf{Phần 2. Trắc nghiệm đúng sai}
\setcounterpageref{ex}{0}

\begin{ex} %[3D2V29]
Xét tính đúng sai của các khẳng định về các phép tính nhân, chia phạm vi 100:
\choiceTF
{\True Một số tự nhiên chẵn khi chia cho 2 luôn luôn không có số dư.}
{Tích của hai số lẻ có thể là một số chẵn.}
{\True Trong biểu thức có cả phép nhân và phép cộng, ta thực hiện nhân trước, cộng sau.}
{Gấp một số lên 2 lần rồi lại chia cho 2 thì số đó tăng lên 4 lần.}
\loigiai{
\begin{itemize}
    \item Ý a) Đúng: Tính chất của số chẵn.
    \item Ý b) Sai: Sai vì lẻ nhân lẻ luôn là số lẻ.
    \item Ý c) Đúng: Thứ tự thực hiện phép tính.
    \item Ý d) Sai: Sai vì giữ nguyên giá trị ban đầu.
\end{itemize}
\begin{itemize}
    \item a) \textbf{Đúng}.
    \item b) \textbf{Sai}.
    \item c) \textbf{Đúng}.
    \item d) \textbf{Sai}.
\end{itemize}
}
\end{ex}

\begin{ex} %[3D2V29]
Bác Ba có 72 quả cam mang ra chợ bán. Buổi sáng bác bán được $\dfrac{1}{3}$ số cam. Buổi chiều bác bán được một nửa số cam còn lại:
\choiceTF
{\True Buổi sáng bác Ba bán được 24 quả cam.}
{\True Sau buổi sáng bác Ba còn lại 48 quả cam.}
{\True Buổi chiều bác Ba bán được 24 quả cam.}
{Sau cả hai buổi bác Ba đã bán hết toàn bộ số cam.}
\loigiai{
\begin{itemize}
    \item Ý a) Đúng: $72 : 3 = 24$ quả.
    \item Ý b) Đúng: $72 - 24 = 48$ quả.
    \item Ý c) Đúng: $48 : 2 = 24$ quả.
    \item Ý d) Sai: Bác còn lại: $48 - 24 = 24$ quả cam.
\end{itemize}
\begin{itemize}
    \item a) \textbf{Đúng}.
    \item b) \textbf{Đúng}.
    \item c) \textbf{Đúng}.
    \item d) \textbf{Sai}.
\end{itemize}
}
\end{ex}

\textbf{Phần 3. Trả lời ngắn}
\setcounterpageref{ex}{0}

\begin{ex} %[3D2V29]
Tính giá trị của biểu thức: $75 : 5 + 18 \times 3$.
\shortans{69}
\loigiai{
\textbf{Đáp số:} $69$\\
$75 : 5 = 15$; $18 \times 3 = 54$; $15 + 54 = 69$.
}
\end{ex}

\begin{ex} %[3D2V29]
Có 4 chuồng vịt, mỗi chuồng có 15 con vịt. Người ta bắt đi 18 con vịt. Hỏi còn lại bao nhiêu con vịt?
\shortans{42}
\loigiai{
\textbf{Đáp số:} $42$\\
Tổng số vịt: $15 \times 4 = 60$ con. Còn lại: $60 - 18 = 42$ con.
}
\end{ex}

\begin{ex} %[3D2V29]
Tìm số tự nhiên $x$, biết: $(x - 8) \times 4 = 64$.
\shortans{24}
\loigiai{
\textbf{Đáp số:} $24$\\
$x - 8 = 64 : 4 = 16 \Rightarrow x = 16 + 8 = 24$.
}
\end{ex}

\begin{ex} %[3D2V29]
Một thanh sắt dài $95\text{ cm}$. Người ta cắt đi 3 đoạn, mỗi đoạn dài $25\text{ cm}$. Hỏi thanh sắt còn lại dài bao nhiêu xăng-ti-mét?
\shortans{20}
\loigiai{
\textbf{Đáp số:} $20$\\
Đã cắt đi: $25 \times 3 = 75\text{ cm}$. Còn lại: $95 - 75 = 20\text{ (cm)}$.
}
\end{ex}

\begin{ex} %[3D2V29]
Một đoàn gồm 38 học sinh đi thuyền qua sông, mỗi thuyền chở được 6 học sinh. Cần ít nhất bao nhiêu thuyền để chở hết các học sinh?
\shortans{7}
\loigiai{
\textbf{Đáp số:} $7$\\
$38 : 6 = 6$ (dư 2). Cần ít nhất $6 + 1 = 7$ thuyền.
}
\end{ex}

\begin{ex} %[3D2V29]
Hiện nay anh 15 tuổi, em 5 tuổi. Hỏi hiện nay tuổi anh gấp mấy lần tuổi em?
\shortans{3}
\loigiai{
\textbf{Đáp số:} $3$\\
$15 : 5 = 3$ lần.
}
\end{ex}
